%% v2_results_ko.tex — 제8장: V2 에이전틱 RAG 평가 결과 (한국어판)

\chapter{V2 에이전틱 RAG: 평가 결과}
\label{chap:v2_results_ko}

본 장은 제~\ref{chap:v2_methodology_ko}장에서 기술한 V2 에이전틱 RAG 파이프라인의 평가 결과를 제시한다. 4개의 실험이 보고된다: E1(\textit{v2\_oracle}), E2(\textit{v2\_oracle\_noloop}), E3(\textit{v2\_oracle\_rankproxy}), E4(\textit{v2\_heuristic}). 모든 실험은 V1 기준선과 동일한 380개의 응답 가능 질의에서 평가된다. 주요 지표는 검색 품질에 대한 Recall@10, MRR, nDCG@10이며, 파이프라인 동작에 대한 검증 통과율과 평균 루프 횟수이다.

\section{실험 설정}
\label{sec:v2_setup_ko}

\paragraph{데이터셋.} V1과 동일하다(제~\ref{chap:methodology}장): 1,206개 청크 SEKD 코퍼스에 대한 380개의 응답 가능 질의이며, 동일한 질의 분할 및 정답 레이블을 사용한다.

\paragraph{평가 절차.} 각 질의는 전체 V2 파이프라인(QAA $\to$ PA $\to$ RA $\to$ VA $\to$ 선택적 재검색 $\to$ AA)을 통과한다. 최종 루프 이후의 검색된 청크 세트는 $k{=}10$에서 표준 IR 지표를 사용하여 정답 청크 세트와 비교된다. 평가 스크립트(\texttt{scripts/run\_v2\_eval.py})가 4개의 실험을 순차적으로 자동화한다.

\paragraph{V1 참조값.} 표~\ref{tab:v2_full_comparison_ko}에는 직접 비교를 위해 5개의 V1 기준선이 포함된다. 이 값들은 표~\ref{tab:retrieval_performance}(제~\ref{chap:results}장)에서 가져온 것이다.

% ─────────────────────────────────────────────────────────────────────────────

\section{전체 검색 성능}
\label{sec:v2_overall_ko}

표~\ref{tab:v2_full_comparison_ko}는 9개 시스템 전체(V1 기준선 5개, V2 실험 4개)에 대한 주요 검색 지표를 보고한다.

\begin{table}[htbp]
\centering
\caption{전체 시스템 비교: V1 기준선 및 V2 에이전틱 실험 ($n{=}380$). 굵은 글씨는 열별 최우수 값.}
\label{tab:v2_full_comparison_ko}
\begin{tabular}{lccccc}
\toprule
\textbf{시스템} & \textbf{Recall@10} & \textbf{MRR} & \textbf{nDCG@10} & \textbf{Hit@10} & \textbf{Recall@1} \\
\midrule
\multicolumn{6}{l}{\textit{V1 기준선 (고정 전략)}} \\
BM25                     & 0.713 & 0.457 & 0.509 & 0.768 & 0.298 \\
Dense (BGE-small)        & 0.740 & 0.501 & 0.548 & 0.787 & 0.346 \\
Hybrid (RRF)             & 0.745 & 0.530 & 0.566 & 0.797 & 0.368 \\
\midrule
\multicolumn{6}{l}{\textit{V1 계획기}} \\
규칙 계획기              & \textbf{0.755} & 0.522 & 0.563 & \textbf{0.808} & 0.358 \\
GPT-5 계획기             & 0.750 & \textbf{0.533} & \textbf{0.570} & 0.803 & 0.368 \\
\midrule
\multicolumn{6}{l}{\textit{V2 에이전틱 RAG (본 연구)}} \\
E1: Oracle + 적응형 + 재검색     & 0.743 & 0.529 & 0.566 & 0.797 & 0.368 \\
E2: Oracle + 적응형, 루프 없음   & 0.743 & 0.529 & 0.566 & 0.797 & 0.368 \\
E3: Oracle + 순위프록시 + 재검색 & 0.743 & 0.529 & 0.566 & 0.797 & 0.368 \\
E4: 휴리스틱 + 적응형 + 재검색  & 0.745 & 0.530 & 0.566 & 0.797 & 0.368 \\
\midrule
\multicolumn{6}{l}{\textit{$\Delta$ 규칙 계획기 대비 (최우수 V1 Recall): E1 $-$ 규칙 계획기 (pp)}} \\
\quad & $-$1.2 & $+$0.7 & $+$0.3 & $-$1.1 & $+$1.0 \\
\bottomrule
\end{tabular}
\end{table}

\FloatBarrier

\paragraph{RQ8 --- V2 에이전틱 검색이 V1 계획기를 능가하는가?}

E1(\textit{v2\_oracle})은 Oracle 질의 분석, 적응형 검증 스코어링, 최대 3회의 재검색 루프를 통해 Recall@10$={0.743}$을 달성한다. 이는 최우수 V1 시스템(규칙 계획기, Recall@10$={0.755}$)보다 1.2~pp \emph{낮다}. V2 파이프라인은 모든 지표에서 Hybrid RRF 기준선과 정확히 일치한다(0.745 / 0.530 / 0.566 / 0.797 / 0.368). V2는 BM25(+3.1~pp Recall@10)와 Dense(+0.3~pp)보다는 우수하여 V1 인프라가 에이전틱 래퍼 내에서 올바르게 작동하고 있음을 확인한다. \textbf{답변: 아니오 --- V2 파이프라인은 최우수 V1 계획기를 능가하지 못한다.} 그 이유는 제~\ref{sec:v2_pipeline_stats_ko}절에서 논의되는 파이프라인 통계에 있다: 검증 에이전트의 임계값이 너무 허용적이어서 재검색 루프가 380개 질의 중 5개에서만 발동된다.

% ─────────────────────────────────────────────────────────────────────────────

\section{제거 연구: V2 구성 요소별 기여}
\label{sec:v2_ablation_ko}

4개의 실험이 구조화된 제거 연구를 형성한다. 표~\ref{tab:v2_ablation_ko}는 세 가지 설계 결정의 기여를 분리한다: 재검색 루프(E1 vs.\ E2), 검증 스코어러(E1 vs.\ E3), QA 방법(E1 vs.\ E4).

\begin{table}[htbp]
\centering
\caption{V2 제거 연구: 실험 구성의 쌍대 비교. $\Delta$는 절대 성능 차이(퍼센트 포인트).}
\label{tab:v2_ablation_ko}
\begin{tabular}{llccc}
\toprule
\textbf{비교} & \textbf{제거 구성 요소} & \textbf{$\Delta$ Recall@10} & \textbf{$\Delta$ MRR} & \textbf{$\Delta$ nDCG@10} \\
\midrule
E1 vs.\ E2 & 재검색 루프 (RQ9)              & $0.0$ & $0.0$ & $0.0$ \\
E1 vs.\ E3 & 적응형 vs.\ 순위프록시 스코어러 (RQ10) & $0.0$ & $0.0$ & $0.0$ \\
E1 vs.\ E4 & Oracle vs.\ 휴리스틱 QA (RQ11) & $-0.1$ & $0.0$ & $-0.1$ \\
\bottomrule
\end{tabular}
\end{table}

\paragraph{RQ9 --- 재검색 루프가 성능을 향상시키는가?}

E1(max\_loops${=}3$)과 E2(max\_loops${=}0$)를 비교하면 \textbf{모든 지표에서 차이가 없다}: 두 실험 모두 Recall@10$={0.743}$, MRR$={0.529}$, nDCG@10$={0.566}$을 달성한다. 재검색 루프는 측정 가능한 개선을 제공하지 않는다. 제~\ref{sec:v2_pipeline_stats_ko}절에 자세히 설명된 바와 같이, E1에서 380개 질의 중 단 5개만이 재검색을 유발하였으며, 5개 모두 검증 에이전트가 최대 루프 횟수(3회)를 소진한 후에도 검색 세트를 승인하지 않았다. \textbf{답변: 아니오 --- 재검색 루프는 이 평가에서 순 개선 효과가 없다.} 핵심 발견은 검증 에이전트의 적응형 스코어러가 실제 recall이 낮은 예외 및 시간적 질의에서도 거의 모든 1차 검색을 승인(98.7\%)한다는 것이다.

\paragraph{RQ10 --- 어떤 검증 스코어러가 가장 효과적인가?}

E1(적응형 스코어러)과 E3(순위-프록시 스코어러)를 비교해도 \textbf{모든 지표에서 차이가 없다}. 특히 E3는 100\%의 검증 통과율(표~\ref{tab:v2_pipeline_stats_ko})을 달성한다: 순위-프록시 스코어러는 어떤 1차 검색도 거부하지 않으므로 루프가 완전히 비활성 상태이다. 적응형 스코어러의 약간 더 엄격한 임계값(98.7\% vs.\ 100\% 통과율)도 거부된 질의에 대한 검색 개선으로 이어지지 않는다. \textbf{답변: 현재 임계값 설정 하에서 두 스코어러 모두 우세하지 않다.}

\paragraph{RQ11 --- Oracle 질의 분석이 얼마나 기여하는가?}

E1(Oracle)과 E4(휴리스틱 QA)를 비교하면 Recall@10와 nDCG@10에서 무시할 수 있는 $-0.1$~pp 차이가 나타나며, E4가 \emph{약간} E1을 능가한다(E4: 0.745 vs.\ E1: 0.743). 이 반직관적 결과는 휴리스틱 QA 에이전트가 대부분의 질의에 대해 기본값으로 Hybrid를 선택하기 때문이다(특정 예외 또는 시간적 패턴이 감지되지 않을 때 Hybrid 기본값이 발동). Hybrid RRF가 가장 강력한 고정 전략 기준선이므로 이 기본값이 유리하게 작용한다. E1의 Oracle 추론 유형은 예외 질의를 Dense로, 시간적 질의를 BM25로 라우팅하지만 검증 임계값이 루프 발동을 막기 때문에 이 Oracle 라우팅 우위는 실현되지 않는다. \textbf{답변: 이 평가에서 Oracle QA는 휴리스틱 QA 대비 사실상 이점을 제공하지 않는다.}

% ─────────────────────────────────────────────────────────────────────────────

\section{추론 유형별 성능}
\label{sec:v2_by_type_ko}

표~\ref{tab:v2_by_type_ko}는 E1(최우수 V2 구성)과 유형별 최우수 V1 기준선의 Recall@10을 추론 유형별로 보고한다. 가장 중요한 비교는 가장 어려운 두 범주인 \texttt{exception}과 \texttt{temporal}이다.

\begin{table}[htbp]
\centering
\caption{추론 유형별 Recall@10: 최우수 V1 시스템 vs. V2 E1 (Oracle + 적응형 + 재검색). $\Delta{=}$E1 빼기 유형별 최우수 V1.}
\label{tab:v2_by_type_ko}
\begin{tabular}{lcccc}
\toprule
\textbf{추론 유형} & \textbf{$n$} & \textbf{최우수 V1} & \textbf{E1 (V2)} & \textbf{$\Delta$ (pp)} \\
\midrule
single\_hop    & 220 & 0.836$^*$ & 0.823 & $-$1.3 \\
multi\_hop     & 53  & 0.538$^*$ & 0.519 & $-$1.9 \\
aggregation    & 35  & 0.924$^*$ & 0.914 & $-$1.0 \\
comparison     & 33  & 0.879$^*$ & 0.848 & $-$3.1 \\
exception      & 34  & 0.412$^*$ & 0.382 & $-$3.0 \\
temporal       & 5   & 0.200$^*$ & 0.200 & $\phantom{-}$0.0 \\
\midrule
\textbf{전체}  & 380 & 0.755 & 0.743 & $-$1.2 \\
\bottomrule
\end{tabular}
{\footnotesize $^*$ 유형별 최우수 V1: aggregation/multi\_hop$=$BM25; comparison/single\_hop$=$규칙 계획기 또는 Hybrid; exception$=$BM25 또는 Hybrid; temporal$=$BM25.}
\end{table}

\FloatBarrier

\paragraph{RQ12 --- 재검색이 예외 및 시간적 질의 실패를 해결하는가?}

유형별 결과는 \textbf{V2가 어떤 추론 유형에서도 V1을 개선하지 못함}을 보여준다. 예외 행이 특히 중요하다: E1은 Recall@10$={0.382}$를 달성하는데, 이는 최우수 V1 결과(BM25 또는 Hybrid, 0.412)보다 3.0~pp \emph{낮다}. 시간적 질의는 Recall@10$={0.200}$으로 유지되어 V1 최우수인 BM25와 동일하다.

그 이유는 파이프라인 통계로 확인된다: 검증 에이전트는 E1에서 예외 및 시간적 질의의 100\%를 1차 검색 통과로 승인한다(표~\ref{tab:v2_pipeline_stats_ko}). Dense 검색에서 예외 질의의 실제 recall이 구조적으로 낮음에도 불구하고, 적응형 스코어러는 검색된 청크에 예외 유형 임계값($\theta=0.22$)을 초과하는 코사인 유사도 점수를 부여한다. 문맥적으로 그럴듯한 정책 청크 --- 예외 조항이 아닌 일반 규칙이더라도 --- 이 임계값을 초과하여 검증을 통과한다. 따라서 루프는 바로 이 루프가 대응하도록 설계된 질의 유형에 대해 한 번도 발동되지 않는다. \textbf{답변: 현재 임계값 설정 하에서 재검색 메커니즘은 예외 또는 시간적 실패 모드를 해결하지 못한다.}

% ─────────────────────────────────────────────────────────────────────────────

\section{파이프라인 통계}
\label{sec:v2_pipeline_stats_ko}

표~\ref{tab:v2_pipeline_stats_ko}는 V2 파이프라인의 내부 동작 통계를 보고한다. 이 지표는 재검색이 얼마나 자주 유발되었는지, 필요한 루프 수는 얼마인지, 그리고 모든 루프 소진 후에도 검증 실패가 발생한 질의의 비율을 보여준다.

\begin{table}[htbp]
\centering
\caption{V2 파이프라인 통계 ($n{=}380$). 검증 통과율은 재검색 없이 검증 에이전트가 검색 세트를 승인한 질의의 비율.}
\label{tab:v2_pipeline_stats_ko}
\begin{tabular}{lccccc}
\toprule
\textbf{실험} & \textbf{검증 통과율} & \textbf{평균 루프} & \textbf{0 루프} & \textbf{1--2 루프} & \textbf{$\geq$3 루프} \\
\midrule
E1: Oracle + 적응형 + 루프      & 0.987 & 0.039 & 375 & 0 & 5 \\
E2: Oracle + 적응형, 루프 없음  & 0.987 & 0.000 & 380 & --- & --- \\
E3: Oracle + 순위프록시 + 루프  & 1.000 & 0.000 & 380 & 0 & 0 \\
E4: 휴리스틱 + 적응형 + 루프   & 0.976 & 0.071 & 371 & 0 & 9 \\
\bottomrule
\end{tabular}
\end{table}

파이프라인 통계는 V2 평가의 핵심 발견을 드러낸다: \textbf{재검색 루프가 사실상 비활성 상태이다.} E1(주요 구성)에서 380개 질의 중 375개(98.7\%)가 1차 검색 시도에서 검증을 통과한다. 검증에 실패한 5개 질의(모두 단일 홉 유형)는 모든 3회 루프를 소진한다. E3에서는 순위-프록시 스코어러가 모든 1차 검색을 통과시켜 루프가 단 한 번도 발동되지 않는다.

더 중요한 것은, 파이프라인 통계(pipeline\_stats.json)의 유형별 세부 분석이 E1에서 예외, 다중 홉, 시간적, 집계 및 비교 질의의 1차 검증 통과율이 100\%임을 보여준다는 점이다. 단일 홉 질의만 97.7\%의 1차 통과율을 보인다. 이 패턴은 루프가 대응하도록 설계된 것과 정반대이다: 어려운 질의 유형(예외, 시간적)이 실제 recall이 낮음에도 검증을 통과하여 재검색이 전혀 유발되지 않는다.

\paragraph{RQ13 --- V2 에이전틱 파이프라인의 주요 실패 모드는 무엇인가?}

표~\ref{tab:v2_failure_modes_ko}는 모든 루프 시도를 소진한 질의들의 실패 이유 분포를 보고한다.

\begin{table}[htbp]
\centering
\caption{검증 실패 이유 분포 (E1, 모든 루프 소진 질의).}
\label{tab:v2_failure_modes_ko}
\begin{tabular}{lcc}
\toprule
\textbf{실패 이유} & \textbf{횟수} & \textbf{실패 중 \%} \\
\midrule
insufficient\_coverage & 5 & 100.0\% \\
wrong\_strategy        & 0 & 0.0\% \\
too\_narrow            & 0 & 0.0\% \\
missing\_entity        & 0 & 0.0\% \\
low\_quality           & 0 & 0.0\% \\
no\_chunks             & 0 & 0.0\% \\
\midrule
\textbf{총 실패} & 5 & 100\% \\
\bottomrule
\end{tabular}
\end{table}

모든 5개의 지속적 실패는 \texttt{insufficient\_coverage}로 분류된다: 검증 에이전트가 청크를 검색하였으나 모든 재검색 시도 후에도 유형별 임계값 이상을 달성한 청크가 $\delta{=}3$개 미만이었다. 예외 및 시간적 질의에 대한 \texttt{wrong\_strategy} 실패가 전혀 없는 것은 루프가 이러한 유형에 전혀 관여하지 않았음을 확인한다: 예외 및 시간적 질의는 모두 실제 recall이 낮음에도 불구하고 초기 검증을 통과하므로 이러한 유형에 대한 실패 진단이 발동되지 않는다.

\textbf{답변(RQ13)}: 주요 실패 모드는 소수의 단일 홉 질의(5개 / 380개 = 1.3\%)에 대한 \texttt{insufficient\_coverage}이다. 그러나 더 중요한 발견은 \emph{에이전틱 아키텍처 자체}의 지배적 실패가 검증 임계값 조정에 있다는 것이다: 강한 신호 질의 유형에 충분히 허용적인 임계값이 약한 신호 질의 유형에도 너무 허용적이어서, 루프가 가장 필요한 곳에서 관여하지 못하게 된다.

% ─────────────────────────────────────────────────────────────────────────────

\section{응답 품질 분석}
\label{sec:v2_answer_quality_ko}

응답 에이전트는 인라인 인용 마커 $[1]\ldots[N]$이 있는 텍스트 응답을 생성한다. 평가 스크립트는 평가 환경에서 LLM API 키가 구성되지 않은 경우 상위 3개 검색 청크를 연결하는 추출적 폴백을 구현한다. 평가가 LLM API 키 없이 실행되었으므로 모든 응답 품질 지표(인용 청크 수, 신뢰도, 증거 커버리지)는 추출적 폴백 모드를 반영하며, 이 값들은 모든 실험에서 균일하다(3.0 인용 청크, 구조적으로 1.0의 커버리지). 응답 에이전트가 제공하도록 설계된 LLM 합성 품질을 반영하지 않으므로 여기서는 보고하지 않는다. LLM 판정 프로토콜을 사용한 응답 품질 별도 평가는 향후 연구 과제로 식별된다.

% ─────────────────────────────────────────────────────────────────────────────

\section{V2 연구 질문 응답 요약}
\label{sec:v2_rq_summary_ko}

\paragraph{RQ8 --- V2가 전체 검색 품질에서 최우수 V1 계획기보다 우수한가?} \textbf{아니오.} E1은 Recall@10$={0.743}$을 달성하여 규칙 계획기(0.755)보다 낮다($-$1.2~pp). V2는 Hybrid RRF 기준선과 일치하지만 어떠한 계획기도 능가하지 못한다.

\paragraph{RQ9 --- 반복 재검색 루프가 기여하는가?} \textbf{아니오.} E1과 E2는 동일한 지표를 산출한다(0.743 / 0.529 / 0.566). 380개 질의 중 5개만이 재검색을 유발하며, 루프 반복으로 수정된 것은 없다.

\paragraph{RQ10 --- 적응형 스코어러가 순위프록시보다 효과적인가?} \textbf{측정 가능한 차이 없음.} E1과 E3는 동일한 지표를 산출한다. 순위-프록시 스코어러는 100\% 통과율을 달성하여 루프가 전혀 발동되지 않는다; 적응형 스코어러는 98.7\% 통과율이지만 실패한 5개 질의는 재검색으로도 수정되지 않는다.

\paragraph{RQ11 --- Oracle QA가 휴리스틱 QA 대비 V2 성능에 얼마나 기여하는가?} \textbf{무시할 수 있는 수준.} E4(휴리스틱)는 Recall@10에서 E1(Oracle)을 0.1~pp \emph{약간} 능가한다. 휴리스틱이 대부분의 질의에 대해 강력한 기준선인 Hybrid로 기본 설정되기 때문이다. Oracle 추론 유형은 검증 임계값이 올바른 유형 정보에 따라 루프를 발동시키지 않아 이점을 제공하지 못한다.

\paragraph{RQ12 --- 타겟 재검색이 예외 처리 및 시간적 질의 실패 모드를 해결하는가?} \textbf{아니오.} 예외 recall은 0.382로 유지되어($-$3.0~pp vs.\ 최우수 V1); 시간적 recall은 0.200으로 변화 없다. 두 유형 모두 실제 recall이 낮음에도 1차 검증 통과율이 100\%이어서 재검색이 발동되지 않는다.

\paragraph{RQ13 --- 모든 재검색 루프 소진 후 V2의 주요 실패 모드는 무엇인가?} \textbf{임계값 조정 실패.} 핵심 실패는 검색 또는 계획 구성 요소가 아니라 검증 스코어러에 있다: 현재 임계값 값($\theta_{\text{exception}}{=}0.22$, $\theta_{\text{default}}{=}0.30$)이 약한 신호 질의 유형에 너무 허용적이어서 루프가 불량 검색을 양호한 것으로 분류한다. 검증에 실패한 5개 질의 중 모두 \texttt{insufficient\_coverage}로 분류된다 --- 루프가 발동될 때 문제를 올바르게 식별하지만 전략 전환만으로는 해결할 수 없음을 나타낸다.
